\documentclass[11pt]{article}
\usepackage{amsmath,amssymb,amsthm}
\usepackage{geometry}
\usepackage{booktabs}
\usepackage{graphicx}
\usepackage{hyperref}
\usepackage{enumerate}
\usepackage{listings}
\usepackage{xcolor}
\usepackage{array}
\usepackage{enumitem}
\usepackage[T1]{fontenc}
\usepackage[utf8]{inputenc}
\geometry{letterpaper, margin=1in}

\lstset{
  basicstyle=\ttfamily\small,
  keywordstyle=\color{blue},
  commentstyle=\color{gray},
  stringstyle=\color{red},
  breaklines=true,
  columns=fullflexible,
  keepspaces=true,
  language=C++
}

\newcommand{\file}[1]{\texttt{#1}}
\newcommand{\type}[1]{\texttt{#1}}
\newcommand{\field}[1]{\texttt{#1}}
\newcommand{\op}[1]{\texttt{#1}}

\title{\textbf{Formation Engine Methodology}\\
\Large Canonical Structural Bridge Architecture\\[4pt]
\large Atomistic Kernel, Coarse-Grained Extension, and LL-FIRE Minimization}
\author{Formation Engine Development Team}
\date{Version 3.0.0 --- Revised Technical Draft}

\begin{document}
\maketitle
\begin{abstract}
This section documents the canonical bridge architecture that separates the
formation engine into three deliberately independent layers: the atomistic
kernel, the canonical structural document layer, and the desktop visualization
workstation. The bridge is not a convenience wrapper. It is the architectural
constraint that allows simulation, persistence, and presentation to evolve
without entangling their internal representations. The coarse-grained
extension preserves these invariants while introducing anisotropic
surface-mapped bead descriptors, environment-responsive coupling, and a
lattice-level FIRE minimizer that operates entirely through the same bridge
contract. The purpose of this document is therefore threefold: first, to
record the atomistic bridge implementation as it exists in the codebase;
second, to make explicit the invariants that preserve extensibility,
reproducibility, and front-end independence; and third, to document the
coarse-grained layer and its LL-FIRE minimizer as a concrete demonstration
that the architecture scales without exception.
\end{abstract}

\tableofcontents
\newpage

% =============================================================================
\section{Introduction}
% =============================================================================

The formation engine is organized around a strict three-layer architecture.
Computational chemistry algorithms, canonical structural data, and
visualization or analysis tools remain separate and communicate only through a
shared structural representation.

\[
  \underbrace{\text{Computational Chemistry}}_{\text{atomistic kernel}}
  \;\longleftrightarrow\;
  \underbrace{\text{Canonical Structure / State Data}}_{\text{bridge layer}}
  \;\longleftrightarrow\;
  \underbrace{\text{Visual / Desktop Output}}_{\text{Qt workstation}}
\]

This separation is not merely stylistic. It is an invariant of the system. No
component outside the computational kernel is permitted to operate directly on
algorithm-specific internal data structures. Every inter-layer exchange must
pass through the canonical representation.

That decision gives the engine a stable internal contract. The physics layer
may change its models, numerical procedures, or internal bookkeeping; the
visualization layer may change its widgets, renderer, or workflow; the file
pipeline may add new readers and writers. None of those changes requires the
other layers to understand the private implementation details of the rest.

The present section documents the architecture as implemented in the current
codebase. It covers the three layers, the bridge implementation, the
conversion pathway, and the design rules that enforce separation. Sections
8--10 extend the treatment to the coarse-grained bead model and its
lattice-level FIRE minimizer, demonstrating that the bridge contract scales
to a fundamentally different physical representation without architectural
compromise.
% =============================================================================
\section{Layer One -- Computational Chemistry}
% =============================================================================

\subsection{The Canonical Simulation State}

The computational layer operates on a single data structure defined in
\file{atomistic/core/state.hpp}. This structure represents the complete
instantaneous configuration of an $N$-particle system:

\begin{lstlisting}
// atomistic/core/state.hpp
struct State {
    uint32_t N{};
    std::vector<Vec3>   X;      // positions      (N x 3, Angstrom)
    std::vector<Vec3>   V;      // velocities     (N x 3, Ang/fs)
    std::vector<double> T;      // per-particle temperature proxy
    std::vector<double> Q;      // charges        (N, elementary charge e)
    std::vector<double> M;      // masses         (N, amu)
    std::vector<uint32_t> type; // species id     (N)

    std::vector<Edge>   B;      // bond topology  (graph edges)
    std::vector<Event>  L;      // event log

    std::vector<Vec3>   F;      // forces         (N x 3, kcal/mol/Ang)
    EnergyTerms         E{};    // energy ledger

    BoxPBC              box;    // periodic boundary conditions
};
\end{lstlisting}

Formally, the state at time $t$ is the tuple
\begin{equation}
  \mathcal{S}(t) = \bigl(N,\; \mathbf{M},\; \boldsymbol{\tau},\;
                         \mathbf{X}(t),\; \mathbf{V}(t),\; \mathbf{Q}(t),\;
                         \mathbf{B}(t),\; \mathbf{F}(t),\; \mathcal{E}(t),\;
                         \mathcal{L},\; \mathcal{B}\bigr)
\end{equation}

where $\mathcal{B}$ encodes the periodic simulation cell via its three edge
lengths $(L_x, L_y, L_z)$ and the minimum-image displacement rule:
\begin{equation}
  \Delta\mathbf{r}_{ij}^{\text{MIC}}
    = \Delta\mathbf{r}_{ij}
    - \mathbf{L} \cdot \operatorname{round}\!\left(
        \frac{\Delta\mathbf{r}_{ij}}{\mathbf{L}}\right)
  \quad (\text{component-wise})
\end{equation}

The energy ledger decomposes the total potential energy into additive terms:
\begin{equation}
  U_{\mathrm{total}} = U_{\mathrm{bond}}
                      + U_{\mathrm{angle}}
                      + U_{\mathrm{tors}}
                      + U_{\mathrm{vdW}}
                      + U_{\mathrm{Coul}}
                      + U_{\mathrm{ext}}
                      + U_{\mathrm{pol}}
\end{equation}

This decomposition is recorded in the \type{EnergyTerms} struct and carried
forward to the bridge layer, where it is stored as named scalar properties on
the output frame. The important point is that the kernel does not export a
special-case visualization object. It exports state.

\subsection{Sanity Invariant}

A helper function \op{sane()} enforces a minimal consistency check on any
\type{State} before it is passed to a kernel operation:

\begin{lstlisting}
inline bool sane(const State& s) {
    if (s.N == 0) return false;
    if (s.X.size()    != s.N) return false;
    if (s.V.size()    != s.N) return false;
    if (s.Q.size()    != s.N) return false;
    if (s.M.size()    != s.N) return false;
    if (s.type.size() != s.N) return false;
    return true;
}
\end{lstlisting}

All vectors must have length $N$. A state with $N = 0$ or mismatched array
lengths is rejected before any kernel operation executes. This check is
simple, but it establishes the baseline contract on which the rest of the
system depends.
\subsection{Kernel Operations}

The computational layer exposes a compact set of algorithms. Each operates on a
\type{State} in place, modifying \field{X}, \field{V}, \field{F}, and
\field{E} as appropriate, and returns a statistics object suitable for the
calling context.

\begin{center}
\begin{tabular}{lll}
\toprule
Algorithm & Implementation file & Output statistics \\
\midrule
Force evaluation (single-point) & \file{atomistic/models/}$*$ & \type{EnergyTerms} \\
FIRE energy minimization & \file{atomistic/integrators/fire.cpp} & \type{FIREStats} \\
Velocity Verlet (NVE) & \file{atomistic/integrators/velocity\_verlet.cpp} & \type{VelocityVerletStats} \\
Langevin dynamics (NVT) & \file{atomistic/integrators/velocity\_verlet.cpp} & \type{LangevinStats} \\
Maxwell--Boltzmann init & \file{atomistic/core/maxwell\_boltzmann.cpp} & --- \\
Bond inference & \file{include/io/xyz\_format.hpp} & bond list \\
\bottomrule
\end{tabular}
\end{center}

The force-model interface is:

\begin{lstlisting}
// atomistic/models/model.hpp
struct IModel {
    virtual ~IModel() = default;
    virtual void eval(State& s, const ModelParams& p) const = 0;
};

std::unique_ptr<IModel> create_lj_coulomb_model();
\end{lstlisting}

Calling \op{eval()} fills \field{s.F} and \field{s.E} for the current
positions \field{s.X}, charges \field{s.Q}, and parameters such as $r_c$ and
$k_{\mathrm{Coul}}$. The model has no knowledge of documents, widgets, files,
or presentation logic. That ignorance is a feature, not a limitation.

\subsection{Structural Parsers and Compilers}

Because the computational layer predates the canonical bridge, conversion
functions are provided in two subordinate namespaces:

\begin{lstlisting}
// atomistic/parsers/xyz_parser.hpp
namespace atomistic::parsers {
    State from_xyz (const vsepr::io::XYZMolecule& mol);
    State from_xyza(const vsepr::io::XYZMolecule& mol);
}

// atomistic/compilers/xyz_compiler.hpp
namespace atomistic::compilers {
    vsepr::io::XYZMolecule to_xyz(
        const State& s,
        const std::vector<std::string>& element_names);

    bool save_xyza(
        const std::string& filename,
        const State& s,
        const std::vector<std::string>& element_names);
}
\end{lstlisting}

These functions handle atomic-mass lookup via a lazy-loaded
\type{vsepr::PeriodicTable} singleton (\file{src/pot/periodic\_db.hpp}),
type-map construction, bond-edge conversion, and optional velocity or charge
extraction. They form the lower half of the conversion pathway described in
Section~\ref{sec:pipeline}.
% =============================================================================
\section{Layer Two -- Canonical Structure and State Data}
% =============================================================================

\subsection{Design Principle}

The canonical data layer provides the representation shared by all major
components of the platform. It serves three purposes:

\begin{enumerate}
  \item to decouple simulation algorithms from visualization and user
        interfaces;
  \item to provide a persistent representation for simulation artifacts; and
  \item to allow multiple front-end tools to interpret the same structural
        data without requiring algorithm-specific logic.
\end{enumerate}

The concrete implementation of this layer is \type{scene::SceneDocument},
defined in \file{apps/desktop/scene/SceneDocument.h}. The desktop never
operates on \type{atomistic::State} directly.

\subsection{SceneDocument}

\begin{lstlisting}
// apps/desktop/scene/SceneDocument.h
struct SceneDocument {
    std::vector<FrameData>                       frames;
    std::map<std::string, PropertyValue>         properties;
    Provenance                                   provenance;

    bool       empty()         const { return frames.empty(); }
    int        frame_count()   const { return (int)frames.size(); }
    bool       is_trajectory() const { return frames.size() > 1; }

    const FrameData& current_frame() const { return frames.back(); }
    FrameData&       current_frame()       { return frames.back(); }
};
\end{lstlisting}

A document holds an ordered sequence of \type{FrameData} snapshots. A
single-frame document represents a static structure. A multi-frame document
represents a trajectory. This distinction is exposed through
\op{is\_trajectory()} and gives the bridge a uniform way to represent both
one-shot results and time-resolved simulations.

\subsection{FrameData}

Each frame is a self-contained snapshot:

\begin{lstlisting}
struct FrameData {
    // Required fields
    std::vector<AtomRecord>        atoms;
    std::vector<BondRecord>        bonds;

    // Optional per-atom fields (size == N when present)
    std::vector<Vec3d>             velocities;  // Ang/fs
    std::vector<Vec3d>             forces;      // kcal/(mol*Ang)
    std::vector<double>            charges;     // elementary charge e

    // Optional periodic cell
    BoxInfo                        box;

    // Frame-level scalar properties (open key-value map)
    std::map<std::string, double>  properties;

    // Provenance
    int         step{-1};
    double      time{-1.0};       // fs, or -1 if not time-resolved
    std::string source_mode;      // "fire", "md_nvt", "import", ...
};
\end{lstlisting}

The \field{properties} map is intentionally open-ended. The engine adapter
writes named scalars such as \texttt{"energy\_total"},
\texttt{"energy\_vdw"}, \texttt{"energy\_coul"},
\texttt{"energy\_bond"}, \texttt{"force\_rms"}, and
\texttt{"temperature"} into this map. The \type{PropertiesPanel} widget then
reads those values by name. In other words, semantics travel as stable keys,
not as privileged knowledge embedded in the UI.

\subsection{AtomRecord and BondRecord}

\begin{lstlisting}
struct AtomRecord {
    int         Z;       // atomic number (identity for rendering)
    Vec3d       pos;     // Cartesian position (Ang)
    std::string symbol;  // element symbol cache ("H", "C", "O", ...)
};

struct BondRecord {
    int    i, j;         // atom indices (0-based)
    double order{1.0};   // 1 = single, 1.5 = aromatic, 2 = double, ...
};
\end{lstlisting}

Atomic number \field{Z} drives the CPK color and van der Waals radius lookup
in \type{ViewportWidget} (Table~\ref{tab:cpk}). The element symbol string
serves as the round-trip carrier across file formats and bridge conversions.

\subsection{Provenance}

\begin{lstlisting}
struct Provenance {
    std::string mode;         // "fire", "md_nvt", "emit", "import", ...
    std::string source_file;  // file path if loaded from disk
    std::string formula;      // chemical formula if known
    std::string hash;         // SHA-256 if available
    std::map<std::string, std::string> parameters;
};
\end{lstlisting}

Every document carries provenance describing how it was produced. The
\type{ObjectTree} widget uses \field{source\_file}, falling back to
\field{mode}, as the display name in the scene hierarchy. The
\type{PropertiesPanel} reads \field{formula} for the identity section.

Provenance matters because the document is not merely a render target. It is a
simulation artifact. Once saved, moved, reopened, or reprocessed, it must
still be able to explain what it is and where it came from.

\subsection{File Format Representation}

Because the canonical layer is the persistent form of structural data, it is
serialized using the following formats:

\begin{center}
\begin{tabular}{lll}
\toprule
Extension & Content & Reader / Writer \\
\midrule
\texttt{.xyz}  & Static atomic geometry & \type{XYZReader} / \type{XYZWriter} \\
\texttt{.xyza} & Trajectory + energy metadata & \type{XYZAReader} / \op{save\_xyza()} \\
\texttt{.xyzc} & Checkpoint with extended system data & \file{src/thermal/xyzc\_format.cpp} \\
\bottomrule
\end{tabular}
\end{center}

All three formats are plain text and human-readable. The standard
\texttt{.xyz} format stores $N$ atom lines of the form
\texttt{<element> <x> <y> <z>}, with coordinates in \AA. The
\texttt{.xyza} format appends an energy-annotated comment line and optional
per-atom velocity and charge columns. The \texttt{.xyzc} checkpoint format
adds velocity, charge, and simulation-cell information sufficient to restart
an MD run.

Bond topology is not natively encoded in the XYZ family. It is either inferred
post-load via \op{XYZReader::detect\_bonds()} using a covalent-radius test
with a $1.2\times$ scale factor, or stored separately in structured manifests.
% =============================================================================
\section{Layer Three -- Visualization and Desktop}
% =============================================================================

\subsection{Overview}

The visualization layer consists of a Qt 6 desktop workstation defined in
\file{apps/desktop/}. It consumes \type{scene::SceneDocument} exclusively and
has no knowledge of \type{atomistic::State}, force models, integrators, or
other kernel types.

The workstation layout follows a familiar CAD-style grammar:

\begin{verbatim}
+-------+-----------------------------+-----------+
| Obj   |                             | Properties|
| Tree  |  GL Viewport (central)      | Inspector |
| dock  |                             | dock      |
+-------+----+------------------------+-----------+
| Console / Log  (bottom dock)                    |
+-------------------------------------------------+
| Status bar                                      |
+-------------------------------------------------+
\end{verbatim}

That arrangement is not just visual convention. It reflects the role of the
front end: inspect the document, invoke operations, and display results. The
Qt layer is an operator console, not a second simulation engine.

\subsection{ViewportWidget}

\file{apps/desktop/ViewportWidget.h} hosts an OpenGL 3.3 Core Profile
renderer. It receives a document via \op{setDocument()} and queries the
current frame on each paint cycle.

\begin{lstlisting}
// ViewportWidget.h
void setFrame   (const scene::FrameData& frame);
void setDocument(std::shared_ptr<scene::SceneDocument> doc);
const scene::FrameData* currentFrame() const;
\end{lstlisting}

Rendering proceeds entirely from \type{FrameData}: atoms are drawn as sphere
imposters whose color and radius are indexed by \field{Z} against an inline
CPK table (Table~\ref{tab:cpk}); bonds are drawn as cylinder pairs. The
viewport performs no force evaluation, bond inference, or element
identification. Those decisions have already been resolved upstream in the
canonical layer.

\begin{table}[ht]
\centering
\caption{Selected CPK color entries from the inline table in
\file{apps/desktop/ViewportWidget.cpp} (index $= Z - 1$).}
\label{tab:cpk}
\begin{tabular}{cllll}
\toprule
$Z$ & Symbol & R & G & B \\
\midrule
1  & H  & 0.90 & 0.90 & 0.90 \\
6  & C  & 0.20 & 0.20 & 0.20 \\
7  & N  & 0.19 & 0.31 & 0.97 \\
8  & O  & 0.90 & 0.05 & 0.05 \\
16 & S  & 1.00 & 1.00 & 0.19 \\
17 & Cl & 0.12 & 0.94 & 0.12 \\
\bottomrule
\end{tabular}
\end{table}

\subsection{ObjectTree}

\file{apps/desktop/ObjectTree.h} is a \type{QTreeWidget} with three fixed
root nodes: \emph{Molecules}, \emph{Crystals}, and \emph{Results}. Items are
added by \type{MainWindow::syncPanels()} after each kernel operation:

\begin{lstlisting}
// From MainWindow::syncPanels()
objectTree_->clear();
QString name = QString::fromStdString(
    doc_->provenance.source_file.empty()
        ? doc_->provenance.mode
        : doc_->provenance.source_file);
objectTree_->addMolecule(name, f.atom_count(), f.bond_count());
\end{lstlisting}

The tree therefore reflects provenance rather than internal engine state. A
document loaded from \texttt{water.xyz} appears as ``water.xyz''; one produced
by FIRE relaxation appears as ``fire''.

\subsection{PropertiesPanel}

\file{apps/desktop/PropertiesPanel.h} is a \type{QScrollArea} containing two
\type{QGroupBox} sections:

\begin{itemize}
  \item \textbf{Identity} -- formula (from \field{provenance.formula}) and
        atom count.
  \item \textbf{Energy} -- total energy in kcal/mol (from
        \field{properties["energy\_total"]}) and RMS force (from
        \field{properties["force\_rms"]}).
\end{itemize}

All values are read from the canonical \type{FrameData} by name. The panel
knows nothing about which kernel produced them, which is precisely why it does
not need to change when new simulation modes are added.

\subsection{ConsolePanel}

\file{apps/desktop/ConsolePanel.h} provides a timestamped log output and a
command-entry line. The \op{commandSubmitted()} signal connects to
\op{MainWindow::onCommand()}, which dispatches text commands such as
\texttt{relax}, \texttt{md}, \texttt{sp}, and \texttt{reset} through the
\type{EngineAdapter} using the same request--result pathway as the toolbar
actions.
% =============================================================================
\section{The Bridge -- EngineAdapter}
\label{sec:adapter}
% =============================================================================

\subsection{Role}

\type{bridge::EngineAdapter} (\file{apps/desktop/bridge/EngineAdapter.h}) is
the only desktop class that instantiates and calls into the atomistic kernel.
It translates between the two representations:

\[
  \text{scene::SceneDocument}
  \;\xrightarrow{\text{doc\_to\_state}}\;
  \text{atomistic::State}
  \;\xrightarrow{\text{kernel op}}\;
  \text{atomistic::State}
  \;\xrightarrow{\text{state\_to\_frame}}\;
  \text{scene::SceneDocument}
\]

The desktop calls only two public methods:

\begin{lstlisting}
// Synchronous (blocks calling thread)
KernelResult run(const KernelRequest& req,
                 ProgressFn progress = nullptr);

// Asynchronous (returns immediately; result via std::future)
std::future<KernelResult> runAsync(const KernelRequest& req,
                                   ProgressFn progress = nullptr);
\end{lstlisting}

This makes \type{EngineAdapter} the architectural choke point. That sounds
harsh because it is. Every path between UI and kernel is forced through one
controlled translation boundary, which is exactly how the layers avoid quietly
bleeding into one another.

\subsection{KernelOp Enumeration}

All operations that the desktop may request are listed in a single
enumeration. New chemistry is added here and nowhere else in the UI:

\begin{lstlisting}
enum class KernelOp {
    LoadXYZ,     // parse file -> SceneDocument
    SaveXYZ,     // write current frame to file
    Relax,       // FIRE energy minimization
    MD_NVE,      // velocity Verlet (constant energy)
    MD_NVT,      // Langevin thermostat (constant temperature)
    SinglePoint, // evaluate energy & forces, no motion
    InferBonds,  // re-detect bonds from covalent radii
    EmitCrystal, // build crystal from preset + supercell
};
\end{lstlisting}

This enumeration is the authoritative contract between the UI and the kernel.
Because the desktop constructs only \type{KernelRequest} values and consumes
only \type{KernelResult} values, the internal implementation of each
operation remains hidden.

\subsection{KernelRequest and KernelResult}

\begin{lstlisting}
struct KernelRequest {
    KernelOp op;

    // Parameters (not all used by every op)
    std::string file_path;              // LoadXYZ, SaveXYZ
    int         max_steps    = 1000;    // Relax, MD_*
    double      dt           = 1.0;     // MD_* timestep (fs)
    double      temperature  = 300.0;   // MD_NVT target T (K)
    double      force_tol    = 1e-4;    // Relax convergence
    std::string preset;                 // EmitCrystal name
    int         supercell[3] = {1,1,1}; // EmitCrystal replication

    std::shared_ptr<scene::SceneDocument> input; // existing data (modify ops)
};

struct KernelResult {
    bool        success{false};
    std::string message;
    std::shared_ptr<scene::SceneDocument> output;
};
\end{lstlisting}

Both structs are plain value carriers containing only standard-library types
and a \type{shared\_ptr} to a document. They carry no kernel-internal
pointers, no raw state references, and no handles to FIRE or Langevin
objects.

\subsection{Internal Implementation -- EngineCore}

The adapter uses a pimpl idiom. The actual kernel objects are hidden in a
private \type{detail::EngineCore} struct defined only inside
\file{apps/desktop/bridge/EngineAdapter.cpp}:

\begin{lstlisting}
namespace detail {
struct EngineCore {
    std::unique_ptr<atomistic::IModel> model;
    atomistic::ModelParams             mp;
    std::mt19937                       rng;

    EngineCore() : rng(std::random_device{}()) {
        model = atomistic::create_lj_coulomb_model();
        mp.rc = 10.0;  // 10 Ang nonbonded cutoff
    }
};
} // namespace detail
\end{lstlisting}

The \type{EngineAdapter} header forward-declares this struct but never exposes
it to the rest of the desktop.
% =============================================================================
\section{Conversion Pipeline}
\label{sec:pipeline}
% =============================================================================

\subsection{SceneDocument $\rightarrow$ atomistic::State}

Before any modifying operation (\texttt{Relax}, \texttt{MD\_NVE},
\texttt{MD\_NVT}, or \texttt{SinglePoint}), the adapter converts the current
\type{SceneDocument} frame into a \type{State}. The procedure is as follows:

\begin{enumerate}
  \item Construct a temporary \type{vsepr::io::XYZMolecule} from
        \field{FrameData::atoms} using \field{symbol} and \field{pos}.
  \item Copy bond records into \field{XYZMolecule::bonds}.
  \item Call \op{atomistic::parsers::from\_xyz(mol)}. This populates masses via
        the \type{vsepr::PeriodicTable} singleton, builds the type map, and
        zero-initializes velocities.
  \item If \field{FrameData::velocities} is non-empty, overwrite
        \field{State::V} with the stored velocity vectors.
\end{enumerate}

This route intentionally passes through the canonical parser rather than
re-implementing its mass-lookup and type-map logic inside the adapter.

\subsection{atomistic::State $\rightarrow$ scene::FrameData}

After any modifying operation, the resulting state is packed back into a
\type{FrameData}:

\begin{enumerate}
  \item Copy positions, velocities, and forces component-wise into
        \field{AtomRecord::pos}, \field{FrameData::velocities}, and
        \field{FrameData::forces}.
  \item Recover \field{AtomRecord::Z} from the element symbol.
  \item Write the energy decomposition into the properties map:
        \texttt{energy\_total}, \texttt{energy\_bond},
        \texttt{energy\_vdw}, and \texttt{energy\_coul}.
  \item Compute and store \texttt{force\_rms} as
        $\sqrt{\frac{1}{N}\sum_i |\mathbf{F}_i|^2}$.
  \item Set \field{source\_mode} to the operation name, for example
        \texttt{"fire"}, \texttt{"md\_nvt"}, or
        \texttt{"single\_point"}.
\end{enumerate}

The resulting \type{FrameData} is appended to a copy of the input document.
That choice preserves the original frame as frame 0 and makes each successful
operation extend the trajectory rather than overwrite history.

\subsection{File Load Pathway}

For \op{KernelOp::LoadXYZ}, the conversion does not pass through
\type{SceneDocument} on the way in:

\[
  \texttt{.xyz} \;\xrightarrow{\text{XYZReader}}\;
  \texttt{XYZMolecule} \;\xrightarrow{\text{detect\_bonds}}\;
  \texttt{XYZMolecule} \;\xrightarrow{\text{from\_xyz}}\;
  \texttt{State} \;\xrightarrow{\text{state\_to\_frame}}\;
  \texttt{SceneDocument}
\]

Bond detection uses \op{XYZReader::detect\_bonds()} with a covalent-radius
scale factor of $1.2$. The resulting bond list is stored in
\field{FrameData::bonds} and appears immediately in the viewport and object
tree.

\subsection{File Save Pathway}

For \op{KernelOp::SaveXYZ}:

\[
  \texttt{SceneDocument} \;\xrightarrow{\text{doc\_to\_state}}\;
  \texttt{State} \;\xrightarrow{\text{compilers::to\_xyz}}\;
  \texttt{XYZMolecule} \;\xrightarrow{\text{XYZWriter}}\;
  \texttt{.xyz}
\]

The element symbol list is extracted from
\field{FrameData::atoms[i].symbol} and passed to \op{to\_xyz()} as the
\field{element\_names} argument.
% =============================================================================
\section{Bidirectional Data Flow}
% =============================================================================

Although the architecture is often drawn as a forward pipeline, the bridge is
fundamentally bidirectional.

\begin{itemize}
  \item \textbf{Forward (kernel $\rightarrow$ desktop):} every kernel
        operation returns a \type{KernelResult} whose \field{output}
        document is passed to \op{ViewportWidget::setDocument()} and
        \op{MainWindow::syncPanels()}. The viewport redraws, and the object
        tree and properties panel update from the same canonical source.

  \item \textbf{Reverse (desktop $\rightarrow$ kernel):} every modifying
        \type{KernelRequest} carries an \field{input} field containing the
        current \type{SceneDocument}. The adapter converts that document back
        into a \type{State}, runs the requested algorithm, and returns a new
        document. From the kernel's perspective, the document is both the
        incoming carrier and the outgoing artifact.
\end{itemize}

The console command dispatcher illustrates the full round trip. The
\op{onCommand()} slot in \type{MainWindow} converts a text command into a
\type{KernelRequest}, calls \op{engine\_.run()}, checks
\type{KernelResult::success}, and forwards the result to
\op{syncPanels()}. No kernel type becomes visible in the UI code at any point.
% =============================================================================
\section{The Coarse-Grained Extension}
\label{sec:cg}
% =============================================================================

The coarse-grained layer introduces a fundamentally different physical
representation---anisotropic surface-mapped beads with environment-responsive
coupling---yet it obeys the same bridge invariants established in Sections
2--7. This is not an accident; it is the test that validates the architecture.

\subsection{Bead State}

Where the atomistic kernel operates on point particles carrying mass, charge,
and element identity, the coarse-grained kernel operates on beads carrying
continuous surface descriptors. The CG state for a system of $N_B$ beads is

\begin{equation}
\label{eq:bead-state}
  \mathcal{S}_{\mathrm{CG}} = \bigl\{
    \mathbf{R}_B,\; \hat{\mathbf{n}}_B,\; \mathcal{D}_B,\; X_B
  \bigr\}_{B=1}^{N_B}
\end{equation}

where $\mathbf{R}_B$ is the bead centre position, $\hat{\mathbf{n}}_B$ is the
primary orientation axis, $\mathcal{D}_B$ is the unified surface descriptor,
and $X_B$ is the environment-responsive state. Each of these is documented
below.

\subsection{Spherical Harmonic Surface Descriptors}

Each bead's anisotropic surface is encoded as a truncated real spherical
harmonic expansion. For a single interaction channel $k$ and angular momentum
cutoff $\ell_{\max}$, the surface function is

\begin{equation}
\label{eq:sh-coeffs}
  \sigma_B^{(k)}(\theta, \phi)
    = \sum_{\ell=0}^{\ell_{\max}} \sum_{m=-\ell}^{\ell}
      c_{\ell m}^{(k,B)} \, Y_\ell^m(\theta, \phi)
\end{equation}

where $Y_\ell^m$ are the real spherical harmonics and $c_{\ell m}^{(k,B)}$
are the expansion coefficients stored in a flat vector of length
$(\ell_{\max}+1)^2$. The implementation is in
\file{coarse\_grain/core/surface\_descriptor.hpp} and
\file{coarse\_grain/core/spherical\_harmonics.hpp}.

\subsection{Unified Descriptor}

The unified descriptor (\file{coarse\_grain/core/unified\_descriptor.hpp})
groups three channels under a single inspectable object:

\begin{center}
\begin{tabular}{lll}
\toprule
Channel $k$ & Physical role & Typical $\ell_{\max}$ \\
\midrule
Steric        & Excluded-volume shape  & 2--4 \\
Electrostatic & Charge anisotropy      & 1--3 \\
Dispersion    & Polarizability surface & 0--2 \\
\bottomrule
\end{tabular}
\end{center}

Each channel carries its own coefficient vector, its own $\ell_{\max}$, and
an \field{active} flag. This is deliberately \emph{not} a tiered model
selector. A bead does not choose between ``low'' and ``high'' resolution; it
carries whatever resolution its coefficients describe, and the interaction
engine reads those coefficients directly. The unified descriptor replaced
an earlier discrete model selector precisely because tier-based branching
introduced silent discontinuities that were invisible to the user.
\subsection{Per-Channel Radial Kernels}

The interaction between two beads is decomposed into per-channel, per-$\ell$
contributions. Each channel has a physically motivated radial kernel
$K_k(\ell, r)$ defined in \file{coarse\_grain/core/channel\_kernels.hpp}:

\begin{align}
\label{eq:kernel-steric}
  K_{\mathrm{steric}}(\ell, r)
    &= \frac{\exp(-\alpha_s \, r)}{1 + \ell} \\[4pt]
\label{eq:kernel-elec}
  K_{\mathrm{elec}}(\ell, r)
    &= \frac{1}{r^{\,\ell+1}} \\[4pt]
\label{eq:kernel-disp}
  K_{\mathrm{disp}}(\ell, r)
    &= \frac{-C_6}{r^{\,6+\ell}}
\end{align}

The steric kernel decays exponentially with a $(1+\ell)$ suppression that
damps higher-order contributions at short range. The electrostatic kernel
follows the standard multipole expansion. The dispersion kernel extends the
$r^{-6}$ London interaction to anisotropic geometry. All three are evaluated
for every $(\ell, r)$ pair during the harmonic-space interaction sum.

\subsection{Interaction Engine}

The full pairwise interaction energy between beads $A$ and $B$ is computed by
the interaction engine (\file{coarse\_grain/models/interaction\_engine.hpp}).
The procedure is:

\begin{enumerate}
  \item Compute the separation vector $\mathbf{r}_{AB}$ and distance
        $r = |\mathbf{r}_{AB}|$.
  \item Compute the relative rotation matrix
        $R = Q_A^{\mathsf{T}} Q_B$ from the bead orientation frames.
  \item Rotate $B$'s SH coefficients into $A$'s local frame using Wigner
        $D$-matrices (\file{coarse\_grain/core/sh\_rotation.hpp}):
        \[
          \tilde{c}_{\ell m}^{(k,B)}
            = \sum_{m'=-\ell}^{\ell} D^{(\ell)}_{mm'}(R)\;
              c_{\ell m'}^{(k,B)}
        \]
  \item Sum the per-channel, per-$(\ell, m)$ contributions:
\end{enumerate}

\begin{equation}
\label{eq:interaction}
  U_{AB}
    = \sum_{k} \lambda_k \sum_{\ell=0}^{\ell_{\max}}
      K_k(\ell, r)
      \sum_{m=-\ell}^{\ell}
        c_{\ell m}^{(k,A)} \, \tilde{c}_{\ell m}^{(k,B)}
\end{equation}

where $\lambda_k$ are per-channel coupling weights. The implementation
returns a fully decomposed \type{InteractionResult} containing the total
energy, per-channel energies, and per-$\ell$ breakdowns. Nothing is hidden.

\subsection{Residual-Driven Adaptive Refinement}

The descriptor residual module
(\file{coarse\_grain/models/descriptor\_residual.hpp}) monitors the truncation
error of each channel's SH expansion and can promote a bead's resolution
when the residual exceeds a configurable threshold. Promotion adds
coefficients for the next angular momentum order. There is no demotion path
in the current implementation; refinement is monotone. The key point is that
this operates on the same unified descriptor---there is no second descriptor
type for ``refined'' beads.
\subsection{Environment State}

The environment-responsive extension adds four per-bead variables grouped as

\begin{equation}
\label{eq:env-state}
  X_B = \bigl\{\rho_B,\; C_B,\; P_{2,B},\; \eta_B\bigr\}
\end{equation}

where $\rho_B$ is the local density (Gaussian-weighted neighbour sum), $C_B$
is the coordination number, and $P_{2,B}$ is the second Legendre polynomial
of the inter-orientation angle, measuring local orientational order:

\[
  P_{2,B} = \frac{1}{|\mathcal{N}_B|}
    \sum_{C \in \mathcal{N}_B}
    \tfrac{1}{2}\bigl(3\cos^2\theta_{BC} - 1\bigr)
\]

The first three are fast observables recomputed at every step. The fourth,
$\eta_B \in [0,1]$, is a slow internal state variable that relaxes toward a
target function via first-order dynamics:

\begin{equation}
\label{eq:eta-relax}
  \eta_B(t + \Delta t)
    = f_B + \bigl(\eta_B(t) - f_B\bigr)\,\exp\!\bigl(-\Delta t / \tau\bigr)
\end{equation}

where $f_B = \alpha\,\hat{\rho}_B + \beta\,\hat{P}_{2,B}$ is the target
function, $\hat{\rho}_B$ and $\hat{P}_{2,B}$ are the normalised observables,
and $\tau$ is the relaxation timescale. The implementation provides both
this exact exponential integrator and a forward-Euler fallback.

\subsection{Environment-Responsive Kernel Modulation}

The environment state modifies the interaction kernels through a linear
coupling. Each channel's kernel is scaled by a modulation factor:

\begin{equation}
\label{eq:modulation}
  K_k(\ell, r;\, \eta_A, \eta_B)
    = K_k(\ell, r) \cdot g_k(\eta_A, \eta_B)
\end{equation}

where

\[
  g_k(\eta_A, \eta_B) = 1 + \gamma_k \, \bar{\eta},
  \qquad
  \bar{\eta} = \tfrac{1}{2}(\eta_A + \eta_B)
\]

and the per-channel coupling constants have the following physical
interpretation:

\begin{center}
\begin{tabular}{llll}
\toprule
Channel $k$ & $\gamma_k$ & Default & Physical effect \\
\midrule
Steric        & $\gamma_{\mathrm{steric}}$ & $+0.2$ & Hardening under compression \\
Electrostatic & $\gamma_{\mathrm{elec}}$   & $-0.1$ & Screening under crowding \\
Dispersion    & $\gamma_{\mathrm{disp}}$   & $+0.5$ & Enhancement under ordering \\
\bottomrule
\end{tabular}
\end{center}

A sign invariant is enforced: $|1 + \gamma_k \bar{\eta}| > 0$ for all
channels, so the kernel never changes sign. The implementation is in
\file{coarse\_grain/models/environment\_coupling.hpp}.
% =============================================================================
\section{Lattice-Level FIRE Minimization}
\label{sec:fire}
% =============================================================================

The lattice-level FIRE (LL-FIRE) minimizer
(\file{coarse\_grain/models/bead\_fire.hpp}) brings energy minimization to
the coarse-grained layer. It operates on bead positions using the full
anisotropic interaction engine for force evaluation, including SH rotation,
per-channel kernels, and optional environment modulation.

\subsection{Total Energy}

The objective function is the sum of all pairwise interaction energies:

\begin{equation}
\label{eq:total-energy}
  U_{\mathrm{total}} = \sum_{A < B} U_{AB}
    = \sum_{A < B} \sum_{k} \lambda_k
      \sum_{\ell, m}
        c_{\ell m}^{(k,A)} \, \tilde{c}_{\ell m}^{(k,B)} \, K_k(\ell, r_{AB})
\end{equation}

When environment coupling is active, each kernel $K_k$ is replaced by its
modulated form from Eq.~\eqref{eq:modulation}.

\subsection{Force Evaluation}

Analytical gradients through the SH rotation and per-$\ell$ kernel chain are
nontrivial. The current implementation uses central finite differences on each
Cartesian component of each bead position:

\begin{equation}
\label{eq:fd-forces}
  F_{B}^{\alpha}
    = -\frac{U(\mathbf{R}_B + \delta\,\hat{e}_\alpha)
           - U(\mathbf{R}_B - \delta\,\hat{e}_\alpha)}{2\delta}
\end{equation}

for $\alpha \in \{x, y, z\}$ and $\delta = 10^{-4}$~\AA\ by default. This
requires $6N_B$ energy evaluations per step, each of which is a full
traversal of the pairwise interaction sum. The cost is intentional: it
guarantees that the forces are exactly consistent with the energy surface,
which matters more during development than raw throughput. Analytical
gradients are a planned replacement once the model stabilizes.

\subsection{FIRE Algorithm}

The FIRE procedure~\cite{Bitzek2006} adapts the following steps at each
iteration $t$:

\begin{enumerate}[label=\arabic*.]
  \item Evaluate forces $\mathbf{F}_B$ and total energy $U$ at current
        positions.
  \item Compute RMS force $F_{\mathrm{rms}} = \sqrt{N_B^{-1} \sum_B |\mathbf{F}_B|^2}$
        and maximum per-bead force $F_{\max}$.
  \item Check convergence: if $F_{\mathrm{rms}} < \varepsilon_F$ or
        $|\Delta U|/N_B < \varepsilon_U$, terminate.
  \item Compute power $P = \sum_B \mathbf{v}_B \cdot \mathbf{F}_B$.
  \item Velocity mixing:
        $\mathbf{v}_B \leftarrow
          (1 - \alpha)\,\mathbf{v}_B
          + \alpha\,|\mathbf{v}|\,\hat{\mathbf{F}}_B$,
        where $|\mathbf{v}| = \sqrt{\sum_B |\mathbf{v}_B|^2}$ and
        $\hat{\mathbf{F}}_B = \mathbf{F}_B / |\mathbf{F}|$.
  \item If $P > 0$ for more than $n_{\min}$ consecutive steps:
        $\Delta t \leftarrow \min(\Delta t \cdot f_{\mathrm{inc}},\, \Delta t_{\max})$
        and $\alpha \leftarrow \alpha \cdot f_\alpha$.
  \item If $P \leq 0$: reset $\Delta t \leftarrow \Delta t \cdot f_{\mathrm{dec}}$,
        $\alpha \leftarrow \alpha_0$, and kick velocities along the force
        direction.
  \item Position update: $\mathbf{R}_B \leftarrow \mathbf{R}_B + \Delta t \cdot \mathbf{v}_B$.
\end{enumerate}

Environment states are optionally refreshed every $n_{\mathrm{env}}$ steps
when environment coupling is active.
\subsection{Parameter Table}

\begin{center}
\begin{tabular}{llll}
\toprule
Parameter & Symbol & Default & Unit \\
\midrule
Initial timestep      & $\Delta t_0$         & 1.0   & fs \\
Velocity mixing       & $\alpha_0$           & 0.1   & --- \\
dt increase factor    & $f_{\mathrm{inc}}$   & 1.1   & --- \\
dt decrease factor    & $f_{\mathrm{dec}}$   & 0.5   & --- \\
$\alpha$ decay        & $f_\alpha$           & 0.99  & --- \\
Min positive steps    & $n_{\min}$           & 5     & --- \\
Maximum timestep      & $\Delta t_{\max}$    & 10.0  & fs \\
Force tolerance       & $\varepsilon_F$      & $10^{-4}$ & kcal/(mol$\cdot$\AA) \\
Energy tolerance      & $\varepsilon_U$      & $10^{-8}$ & kcal/mol \\
Maximum steps         & ---                  & 2000  & --- \\
FD step size          & $\delta$             & $10^{-4}$ & \AA \\
\bottomrule
\end{tabular}
\end{center}

All parameters are stored in the inspectable \type{BeadFIREParams} struct.
None are compiled in as magic numbers; all are overridable from the CLI.

\subsection{Convergence Diagnostics}

The minimizer records a \type{BeadFIREStep} struct at every iteration,
containing $U_{\mathrm{total}}$, per-channel energies ($U_{\mathrm{steric}}$,
$U_{\mathrm{elec}}$, $U_{\mathrm{disp}}$), $F_{\mathrm{rms}}$,
$F_{\max}$, $\alpha$, $\Delta t$, and per-bead energy change
$|\Delta U|/N_B$. The full history is stored in
\type{BeadFIREResult::history} and is available for plotting, regression
testing, or post-mortem analysis. This is the anti-black-box philosophy
applied to optimization: every step is a public record.
% =============================================================================
\section{CLI Front-End}
\label{sec:cli}
% =============================================================================

The coarse-grained engine is exposed to the user through a scientific operator
console implemented in \file{src/cli/cg\_commands.cpp}. The CLI serves the
same role for the CG layer that the Qt desktop serves for the atomistic layer:
it consumes canonical state, invokes kernel operations, and displays results.
It does not contain physics.

The command grammar is:

\begin{center}
\begin{tabular}{lp{8cm}}
\toprule
Command & Description \\
\midrule
\texttt{scene}    & Build a bead scene from geometric presets
                    (isolated, pair, stack, T-shape, square, shell, cloud) \\
\texttt{inspect}  & Display bead positions, environment state,
                    unified descriptors \\
\texttt{env}      & Run the environment update pipeline ($\eta$ relaxation)
                    with per-step trajectory reporting \\
\texttt{interact} & Evaluate pairwise interactions with full per-channel,
                    per-$\ell$ decomposition \\
\texttt{viz}      & Launch lightweight GLFW/ImGui bead viewer
                    (when built with \texttt{-DBUILD\_VIS=ON}) \\
\texttt{fire}     & Run LL-FIRE minimization with convergence reporting \\
\bottomrule
\end{tabular}
\end{center}

Each command follows the same pattern: build or load a
\type{CGSystemState}, optionally run environment updates, execute the
requested operation, and print a fully decomposed report. The
\texttt{fire} command, for example, builds a scene, populates unified
descriptors, calls \op{BeadFIRE::minimize()}, and reports the full
convergence trajectory including per-channel energy decomposition at every
reported step.

All commands route through \type{CGSystemState}
(\file{include/cli/system\_state.hpp}), which is the CG analogue of the
atomistic \type{State}. The same principle applies: the CLI never
instantiates kernel objects directly. It operates on the interpretation layer,
and the interpretation layer calls the engine.
% =============================================================================
\section{Architectural Invariants}
% =============================================================================

The canonical bridge architecture enforces five invariants by construction.
Invariants A--C govern the atomistic layer; D and E extend the contract to the
coarse-grained extension.

\paragraph{Invariant A -- Kernel opacity.}
No class in \file{apps/desktop/} except
\file{apps/desktop/bridge/EngineAdapter.cpp} may include headers from
\file{atomistic/}. Types such as \type{atomistic::State},
\type{atomistic::IModel}, \type{FIREParams}, and \type{LangevinParams}
remain invisible to the Qt layer.

\paragraph{Invariant B -- Document as sole inter-layer currency.}
\type{KernelRequest::input} and \type{KernelResult::output} are both
\type{shared\_ptr<scene::SceneDocument>}. Raw simulation arrays, force-model
references, and integrator handles are never passed across the boundary.

\paragraph{Invariant C -- Automatic compatibility of new models.}
Any new \type{KernelOp} that produces a valid \type{scene::SceneDocument}
becomes immediately compatible with the viewport, object tree, properties
panel, console, and file-save pipeline without modification to those
components. Adding a new simulation mode therefore requires only:

\begin{enumerate}
  \item adding an entry to the \type{KernelOp} enum,
  \item adding a case to the \op{run()} dispatch switch in
        \file{EngineAdapter.cpp}, and
  \item optionally adding a UI action in \file{MainWindow.cpp}.
\end{enumerate}

The visualization system itself does not need to be rewritten. That is the
practical payoff of the bridge: once the output satisfies the canonical
contract, the rest of the stack already knows how to consume it.

\paragraph{Invariant D -- Unified descriptor identity.}
A bead does not select between discrete model tiers. It carries a single
\type{UnifiedDescriptor} with three channels, each at its own angular
momentum resolution. The interaction engine reads those channels directly.
There is no switch statement on a tier label, no model-selector dispatch, and
no silent fallback to a simpler approximation. The unified descriptor
replaced an earlier \type{ModelSelector} precisely because tier-based
branching created discontinuities that the bridge could not represent honestly.

\paragraph{Invariant E -- Full inspectability.}
Every intermediate quantity in the CG pipeline---per-$\ell$ kernel values,
per-channel energies, environment observables, modulation factors, FIRE
per-step diagnostics---is stored in a named, typed field and is accessible
from the CLI or any front-end that reads the canonical state. There are no
internal-only diagnostic modes and no quantities that require a debugger to
observe. If a value affects the physics, it appears in the output.
% =============================================================================
\section{Advantages of the Bridge Architecture}
% =============================================================================

\paragraph{Decoupling of Components.}
Simulation algorithms, visualization tools, and analysis modules remain
independent because each communicates only through the canonical structural
representation. Replacing the Qt desktop with a different front end would
require only a new consumer of \type{scene::SceneDocument}, not a rewrite of
the physics code.

\paragraph{Reproducibility.}
Simulation outputs are recorded as explicit structural artifacts. A saved
\texttt{.xyz} or \texttt{.xyza} file can be reloaded, visualized, and
re-simulated independently of the original runtime environment.
\type{Provenance} metadata records operation mode and source information,
making the history of a structure traceable rather than anecdotal.

\paragraph{Multiple Front-End Interfaces.}
Because the visualization layer consumes canonical structural data rather than
internal engine state, multiple front ends can coexist. The current codebase
includes a lightweight ImGui viewer (\file{apps/simple-viewer.cpp}),
the Qt desktop workstation (\file{apps/desktop/}), and the CG scientific
console (\file{src/cli/cg\_commands.cpp}), each consuming canonical state
through its own rendering or reporting path.

\paragraph{Testability.}
The bridge layer is independently testable. A unit test can construct a
synthetic \type{SceneDocument}, submit it as a \type{KernelRequest}, and
assert properties of the returned \type{KernelResult} without launching a Qt
application or opening a file. The CG CLI commands are testable in the same
way: build a \type{CGSystemState}, invoke a kernel operation, and inspect the
result.

\paragraph{Thread Safety.}
\op{EngineAdapter::runAsync()} launches kernel operations on a worker thread
via \op{std::async}. \type{KernelRequest} and \type{KernelResult} are value
carriers copied into and out of that worker context, so no shared mutable
kernel state crosses the thread boundary. The returned
\type{shared\_ptr<scene::SceneDocument>} is then consumed as an output
artifact, not as a live simulation handle.
% =============================================================================
\section{Summary}
% =============================================================================

The canonical structural bridge architecture organizes the formation engine
into three independent layers connected by a single shared data
representation. The coarse-grained extension preserves the same invariants
while introducing anisotropic descriptors, environment coupling, and a
lattice-level energy minimizer.

\begin{center}
\begin{tabular}{lll}
\toprule
Layer & Primary type & Location \\
\midrule
Atomistic kernel     & \type{atomistic::State}       & \file{atomistic/} \\
CG kernel            & \type{BeadSystem}, \type{UnifiedDescriptor} & \file{coarse\_grain/} \\
Canonical bridge     & \type{scene::SceneDocument}    & \file{apps/desktop/scene/} \\
Bridge adapter       & \type{bridge::EngineAdapter}   & \file{apps/desktop/bridge/} \\
CG interpretation    & \type{CGSystemState}           & \file{include/cli/system\_state.hpp} \\
Visualization        & \type{ViewportWidget}, docks   & \file{apps/desktop/} \\
CG CLI               & \op{cg\_dispatch()}            & \file{src/cli/cg\_commands.cpp} \\
LL-FIRE minimizer    & \type{BeadFIRE}                & \file{coarse\_grain/models/bead\_fire.hpp} \\
Interaction engine   & \op{interaction\_energy()}     & \file{coarse\_grain/models/interaction\_engine.hpp} \\
\bottomrule
\end{tabular}
\end{center}

The governing rule is absolute: no component outside the computational kernel
may rely on algorithm-specific internal data structures. All communication
must pass through the canonical representation. That restriction prevents
coupling between simulation physics and interface implementation, preserves
extensibility, and keeps the system intelligible as additional physical models
are introduced.

The coarse-grained extension is the proof. It introduces spherical harmonic
descriptors, Wigner $D$-matrix rotation, per-channel radial kernels,
environment-responsive modulation, and a full FIRE minimizer, and none of
that required the bridge contract to bend. The architecture did not need an
exception. It needed the discipline to not ask for one.

% =============================================================================
\section{Canonical Vector Type: \type{vsepr::Vec3}}
% =============================================================================

\subsection{Motivation}

Every layer in the bridge architecture --- atomistic kernel, bead system,
energy models, XYZ serialization, sensor subsystem, EHD multiphysics, and
analysis pipelines --- requires a three-dimensional position or vector
quantity. Prior to Day~\#56 of the development log, each subsystem carried
its own local \type{Vec3} definition. These definitions were structurally
equivalent but semantically isolated: the atomistic kernel, the UFX schema
layer, the peptide formation engine, the sensor module, and the EHD continuum
solver each declared independent \texttt{struct Vec3} types. The result was a
fragmented vector ecosystem that violated the bridge invariant stated
throughout this document: all inter-layer exchange must pass through a shared,
canonical representation.

The \type{vsepr::Vec3} unification resolves this by replacing every local
definition with a single project-wide type. Compatibility shims were
introduced where callers depended on legacy method names (\texttt{norm\_sq()}
for EHD code; float-space \type{Vec3f} retained in the crystal data model
where single-precision storage is intentional). All other subsystems now
alias directly to the authoritative type.

\subsection{Implementation}

The authoritative definition lives at
\file{src/core/math\_vec3.hpp} and is exposed through the stable canonical
include path \file{src/core/math/vec3.hpp}. The struct is a
plain-data, double-precision aggregate with no external dependencies:

\begin{lstlisting}
namespace vsepr {

struct Vec3 {
    double x{}, y{}, z{};

    // Future interoperability adapter (non-mandatory)
    Eigen::Vector3d to_eigen() const {
        return Eigen::Vector3d{x, y, z};
    }

    static Vec3 from_eigen(const Eigen::Vector3d& v) {
        return {v.x(), v.y(), v.z()};
    }
};

} // namespace vsepr
\end{lstlisting}

The Eigen adapter methods (\texttt{to\_eigen}, \texttt{from\_eigen}) are
present as a migration-safe boundary. They allow individual subsystems to
interoperate with Eigen-backed math at call sites without requiring the rest
of the project to take an Eigen dependency. The adapter is a seam, not an
entanglement.

All previously fragmented definitions are now aliases:

\begin{center}
\begin{tabular}{ll}
\toprule
Subsystem & Alias form \\
\midrule
\file{atomistic/core/state.hpp}          & \texttt{using Vec3 = vsepr::Vec3;} \\
\file{src/v4/uff/ufx\_material\_record.hpp} & \texttt{using Vec3 = vsepr::Vec3;} \\
\file{src/io/xyz\_unified.hpp}           & \texttt{using XYZVec3 = vsepr::Vec3;} \\
\file{chem/peptide/peptide\_formation.hpp} & \texttt{using Vec3 = vsepr::Vec3;} \\
\file{include/sensor/sensor.hpp}         & \texttt{using Vec3 = vsepr::Vec3;} \\
\file{sim/ehd/ehd\_types.hpp}            & \texttt{using Vec3 = vsepr::Vec3;} \\
\bottomrule
\end{tabular}
\end{center}

The aliasing contract is enforced by a compile-time audit test at
\file{tests/vec3\_type\_audit.cpp} (CMake target \texttt{Vec3TypeAudit},
Group~19). The test uses \texttt{static\_assert} to confirm that
\type{atomistic::Vec3} and \type{vsepr::ufx::Vec3} are the same type as
\type{vsepr::Vec3}, and runs smoke tests for \texttt{OnlineStats},
\texttt{StationarityGate}, and core vector operations. The test must pass
before any refactor that touches vector types is considered complete.

\subsection{Design Rationale}

\begin{center}
\begin{tabular}{p{0.32\textwidth}p{0.60\textwidth}}
\toprule
Benefit & Why it matters \\
\midrule
One project-wide type &
  Atomistic, bead, energy, XYZ, and analysis layers all speak the same vector
  language. No conversion overhead, no silent type mismatch between layers. \\[4pt]
Stable file compatibility &
  The XYZ and \texttt{xyzFull} writers know exactly how
  \texttt{position.x}/\texttt{y}/\texttt{z} are stored and serialized.
  The canonical type is the serialization truth. \\[4pt]
No dependency lock-in &
  The project is not tied to Eigen, GLM, Qt, VTK, or any other library's
  vector assumptions. External library vectors are converted at the boundary,
  not threaded through every subsystem. \\[4pt]
Physics-specific helpers &
  Only the helpers VSEPR-SIM actually needs are present: \texttt{norm()},
  \texttt{dot()}, \texttt{cross()}, unit conversions, tolerances, and
  planned periodic boundary condition helpers. No excess API surface. \\[4pt]
Safer audits &
  If every module uses \type{vsepr::Vec3}, duplicate or divergent vector
  definitions become immediately visible to the compiler and to the audit
  test. Regressions cannot hide behind a locally identical struct. \\[4pt]
Cleaner future migration &
  If SIMD acceleration or Eigen-backed math is introduced later, the
  internals of \type{vsepr::Vec3} or its adapters change. No subsystem
  needs to be rewritten because no subsystem holds a private vector type. \\
\bottomrule
\end{tabular}
\end{center}

\subsection{Exception: Float-Space Crystal Data}

The crystal data model (\file{include/data/Crystal.hpp}) retains a separate
\type{Vec3f} (single-precision float) because the crystal lattice and
visualization geometry pipeline deliberately operates in float-space for
memory and rendering performance. This is not a violation of the bridge
invariant; it is a documented precision boundary. The canonical double-precision
\type{vsepr::Vec3} is used at all physics and analysis call sites.
\type{Vec3f} is used only inside the visualization geometry subsystem and
is never exchanged across the bridge.

\begin{thebibliography}{9}
\bibitem{Bitzek2006}
  E.~Bitzek, P.~Koskinen, F.~G\"ahler, M.~Moseler, P.~Gumbsch,
  ``Structural Relaxation Made Simple,''
  \emph{Phys.\ Rev.\ Lett.}, \textbf{97}, 170201 (2006).
\end{thebibliography}

\end{document}